% Notes of the Montreal seminar, 1970 (cote 129, typed part): archive cote 129, site page 67 (PDF page 68).
% typescript. Draft compared: none.
% AI-transcribed by Claude Opus 5.5 (claude-opus-5-5, adaptive thinking, effort high), Message Batches API, batch mtl129_full (msgbatch_01AGnDKHwy6ATCM7ehUa6263), 2026-11-13; not yet checked by eye.
\scanpage{129/67}
\mspage{\hand{III-12}}

$D^*(G)$ est un module sur $W = \underset{\leftarrow}{\lim}\, W_n$; comme limite p-adique
de $W_n$-modules plats, il est plat sur l'anneau local $W$ et est donc
\uline{libre}. Les $M(n)$ étant de longueur finie, $D^*(G)$ est de \uline{type fini}.
De plus, $D^*(G)$ est muni de morphismes semi-linéaires $F$ et $V$ induits
par les morphismes correspondants sur $M(n)$.

6. \uline{Exemples et types particuliers de groupes de Barsotti-Tate}

6.1. Soit $A$ un schéma abélien sur $S$ (i.e. un schéma en groupes
commutatifs, propre et lisse sur $S$, et à fibres connexes). Un schéma
abélien est p-divisible: c'est bien connu dans le cas où la base est
le spectre d'un corps algébriquement clos (cf. S. Lang, ``Abelian
Varieties'') et, si $S$ est quelconque, la surjectivité de $\mathrm{pId}_A$ résulte
de ce cas particulier grâce au critère de platitude fibre par fibre
(EGA, IV, 11.3.10).

Si $d$ désigne la dimension relative de $A$ sur $S$, on peut
montrer que
\[ \text{rang } A(1) = p^{2d} . \]

Le groupe $A(1)$ est donc fini; il est évidemment de présen-%
tation finie et, étant plat, il est localement libre. On en déduit que
$A(\infty) = \underset{\overrightarrow{n}}{\lim}\, A(n)$ est un groupe de BT de rang $2d$. C'est le \uline{p-groupe
de BT associé au schéma abélien $A$}, et cette construction peut se faire
pour tout nombre premier $p$.

Si toutes les caractéristiques résiduelles de $S$ sont diffé-%

\margin{\hand{67}}
